\chapter{Рекомендаційні Системи.}

Рекомендаційна система  - це система, яка автоматично аналізує інформацію про користувачів певного сервісу та об'єкти, якими вони користуються (товари, фільми, ін.) і пропонує користувачам ті об'єкти, в яких вони можуть бути зацікавленими. Тобто, ця система може передбачити, який рейтинг поставить той чи інший користувач певному об'єкту на основі вподобань інших користувачів та вподобань самого цього користувача. 

Тут можна згадати приклад Youtube (алгоритми Youtube самі пропонують користувачам відео на основі їхніх вподобань), інтернет-магазини, Netflix та інші стрімінгові платформи, книжкові сервіси та багато інших. Всі ці онлайн-платформи імплементують алгоритми рекомендаційних систем для аналізу даних, і саме про ці алгоритми ми і поговоримо далі. 

\textit{Контентно-орієнтована фільтрація}: це метод рекомендації, який аналізує характеристики самих об'єктів. Основна ідея полягає у тому, що якщо користувачеві подобалися певні об'єкти, то йому, ймовірно, сподобаються інші об'єкти зі схожими характеристиками (жаер фільму, автор книги, інш.). Серед недоліків цього підходу можна виділити те, що необхідно детально описати кожен об'єкт. 

\section{Колаборативна фільтрація.}

Колаборативна фільтрація (Collaborative Filtering) — це підхід у системах рекомендацій, який прогнозує,
які об’єкти (товари, фільми, книги тощо) можуть сподобатися користувачу,
виходячи з подібності між користувачами або між самими об’єктами. 

Існує декілька підходів до обчислення оцінки:

\textit{Підходи "користувач – користувач" (User-Based Collaborative Filtering)}: знайти користувачів, які найбільш схожі на мене (на основі 
тільки тих елементів, які оцінені для нас обох), і передбачити оцінки для інших 
елементів на основі середнього значення 

\textit{Підходи "елемент – елемент (Item-Based Collaborative Filtering)"}: знайти елементи, найбільш схожі на даний елемент (на основі 
всіх користувачів, які оцінили обидва елементи), і передбачити оцінки для інших користувачів на основі 
середнього значення

\textit{Підходи матричної факторизації}: знайти деяку низькорангову декомпозицію матриці $X$, яка узгоджується з спостережуваними значеннями


Маємо:

\begin{itemize}
\item множину користувачів $ U = {u_1, u_2, \ldots, u_m} $,
\item множину об’єктів $I = {i_1, i_2, \ldots, i_n} $.
\end{itemize}

Створюється матриця оцінок:

$$
R = [r_{ui}]_{m \times n}
$$

де $r_{ui}$ — рейтинг, який користувач $ u $ поставив об’єкту $ i $, якщо користувач не оцінив об’єкт — $ r_{ui} $ відсутній (невідомий).


Передбачити невідомий рейтинг:
$$
\hat{r_{ui}} \approx r_{ui}
$$

для пари (користувач, об’єкт), якої ще немає у матриці.

\subsection{Фільтрація на основі подібності користувачів.}

Розглянемо підхід на основі подібності користувачів (User-based Collaborative Filtering)

Основна ідея: Шукаємо користувачів, схожих на користувача $u $, і прогнозуємо його рейтинг
на основі оцінок цих схожих користувачів.

Кроки:

1. Обчислити схожість між користувачами, наприклад, за допомогою косинусної подібності:
   $$
   \text{sim}(u, v) = \frac{\sum_{i \in I_{uv}} r_{ui} r_{vi}}{\sqrt{\sum_{i \in I_{uv}} r_{ui}^2} \sqrt{\sum_{i \in I_{uv}} r_{vi}^2}}
   $$
   де $ I_{uv} $ — множина об’єктів, які оцінили обидва користувачі.

2. Прогноз рейтингу:
   $$
   \hat{r_{ui}} = \bar{r_u} + \frac{\sum_{v \in N(u)} \text{sim}(u, v) (r_{vi} - \bar{r}_v)}{\sum_{v \in N(u)} |\text{sim}(u, v)|}
   $$
   
де $ N(u) $ — набір найбільш схожих користувачів (сусідів), $\bar{r}_u $ — середній рейтинг користувача $u$.



Окремо можна винести Кореляцію Пірсона, яка дещо за формулою відрізняється від косинусної подібності:
  $$
  \text{sim}(x, y) = \frac{\sum_i (x_i - \bar{x})(y_i - \bar{y})}{\sqrt{\sum_i (x_i - \bar{x})^2} \sqrt{\sum_i (y_i - \bar{y})^2}}
  $$

\begin{table}[htbp]
\centering
\caption{Матриця оцінок користувачів та фільмів}
\label{tab:user_movie_ratings}
\begin{tabular}{ccccc}
\toprule
\textbf{Користувач / Фільм} & \textbf{Фільм A} & \textbf{Фільм B} & \textbf{Фільм C} & \textbf{Фільм D} \\ 
\midrule
\textbf{U1} & 5 & 4 & ? & 2 \\ 
\textbf{U2} & 3 & ? & 4 & 1 \\ 
\textbf{U3} & ? & 2 & 5 & 4 \\ 
\textbf{U4} & 4 & 3 & ? & ? \\ 
\bottomrule
\end{tabular}
\end{table}

На практиці кожен користувач взаємодіє лише з дуже невеликою частиною всіх доступних об'єктів. Наприклад, онлайн-магазин може містити мільйони товарів, але конкретний покупець придбає лише кілька десятків із них. Аналогічно користувач Netflix перегляне лише незначну частину величезної бібліотеки фільмів. Саме тому, на практиці, ця матриця переважно є розрідженою. Ви можете нагадати, як працювати з розрідженими матрицями, в розділі 4.

Запишемо ці дані в numpy масив \verb|ratings_matrix| та перетворимо цю матрицю в датафрейм:

\begin{lstlisting}
import numpy as np
import pandas as pd

ratings_matrix = np.array([
    [5, 3, 0, 1],
    [4, 0, 0, 1],
    [1, 1, 0, 5],
    [1, 0, 0, 4],
    [0, 1, 5, 4],
])

ratings_df = pd.DataFrame(ratings_matrix, columns=["Item1", "Item2", "Item3", "Item4"])
\end{lstlisting}

Обчислимо кореляційну матрицю Пірсона:

\begin{lstlisting}
user_similarity = ratings_df.T.corr(method='pearson')

print("User-User Pearson Correlation Matrix:\n", user_similarity)
\end{lstlisting}

Результат:

\begin{console}
User-User Pearson Correlation Matrix:
           0         1         2         3         4
0  1.000000  0.774291 -0.186441 -0.178683 -0.978839
1  0.774291  1.000000  0.019854  0.162791 -0.628768
2 -0.186441  0.019854  1.000000  0.972828  0.221028
3 -0.178683  0.162791  0.972828  1.000000  0.258904
4 -0.978839 -0.628768  0.221028  0.258904  1.000000
\end{console}


Передбачимо рейтинг користувача для якогось із фільмів:


\begin{lstlisting}
def predict_rating(user_index, item_index, ratings_df, similarity_matrix):
    user_means = ratings_df.replace(0, np.nan).mean(axis=1).values

    item_ratings = ratings_df.iloc[:, item_index]
    
    user_similarities = similarity_matrix[user_index]
    
    non_zero_indices = item_ratings[item_ratings > 0].index
    similarities = user_similarities[non_zero_indices]
    ratings = item_ratings[non_zero_indices]
    means = user_means[non_zero_indices]
    
    if len(similarities) == 0:
        return 0
    
    ratings_weighted = ratings - means
    weighted_sum = np.dot(similarities, ratings_weighted)
    
    sum_of_similarities = np.sum(np.abs(similarities))

    predicted = weighted_sum / sum_of_similarities if sum_of_similarities != 0 else 0

    return ratings.mean() + predicted

print(predict_rating(4, 0, ratings_df, user_similarity))
\end{lstlisting}

Тут ми спочатку отримуємо рейтинг за айтем від усіх інших користувачів:

\begin{lstlisting}
item_ratings = ratings_df.iloc[:, item_index]
\end{lstlisting}

Отримуємо подібність користувача з іншими користувачами:

\begin{lstlisting}
user_similarities = similarity_matrix[user_index]
\end{lstlisting}

Фільтруємо користувачів за тими, хто надав рейтинг цьому айтему і вилучаємо самоподібність:

\begin{lstlisting}
non_zero_indices = item_ratings[item_ratings > 0].index
similarities = user_similarities[non_zero_indices]
ratings = item_ratings[non_zero_indices]
means = user_means[non_zero_indices]
\end{lstlisting}

Якщо жоден інший користувач надав рейтинг цьому айтему, ставимо 0Ж

\begin{lstlisting}
if len(similarities) == 0:
    return 0
\end{lstlisting}

І зрештою, обчислюємо зважену суму рейтингів, нормалізовану за подібностями користувачів:

\begin{lstlisting}
ratings_weighted = ratings - means
weighted_sum = np.dot(similarities, ratings_weighted)
    
sum_of_similarities = np.sum(np.abs(similarities))

predicted = weighted_sum / sum_of_similarities if sum_of_similarities != 0 else 0
\end{lstlisting}

Результат:

\begin{console}
1.0331985858244794
\end{console}

\ding{43} На практиці рекомендаційні системи поєднують декілька методів побудови рекомендацій, таким чином утворюючи гібридну рекомендаційну систему (Hybrid Recommendation System). Саме гібридний підхід використовують більшість сучасних платформ, таких як Amazon і YouTube.

Гібридні системи мають низку переваг: вища точність рекомендацій, часткове розв'язання проблеми холодного старту; менша чутливість до розрідженості даних; більша різноманітність рекомендованих об'єктів; можливість використовувати різні джерела інформації одночасно.Завдяки цьому гібридні системи вважаються сучасним стандартом побудови рекомендаційних сервісів.

\subsection{Холодний старт.}

Проблема холодного старту (Cold Start) виникає тоді, коли рекомендаційна система не має достатньо інформації для формування якісних рекомендацій. Це може бути або коли у нас новий користувач (система іще не знає його вподобань), або новий об'єкт (товар/фільм/книга іще не мають оцінок). Може бути також зовсім нова система вцілому. Через нестачу даних алгоритми колаборативної фільтрації працюють менш ефективно.

Для подолання проблеми холодного старту використовують кілька підходів.

Найпоширенішими є:

\begin{itemize}
\item запропонувати новому користувачеві оцінити декілька об'єктів під час реєстрації;
\item використовувати інформацію про характеристики об'єктів (контентно-орієнтована фільтрація);
\item застосовувати демографічні дані користувачів;
\item рекомендувати популярні або трендові об'єкти;
\item використовувати гібридні рекомендаційні системи;
\item аналізувати поведінку користувача (кліки, перегляди, пошукові запити) навіть без явних оцінок.
\end{itemize}

Завдяки поєднанню цих методів сучасні системи значно зменшують негативний вплив холодного старту.

\section{Матрична факторизація.}





\section{Метрики рекомендаційних систем.}

Для оцінювання рекомендаційних систем використовують різні метрики. Поговоримо про деякі з них.

Precision (точність), як і в задачах класифікації із машинного навчання, показує, яка частина рекомендованих об'єктів дійсно виявилася корисною для користувача.

$$Precision = \frac{TP}{TP+FP}$$

де $TP$ - релевантні рекомендовані об'єкти, $FP$ - нерелевантні.

Recall показує, яку частину всіх релевантних об'єктів система змогла рекомендувати. 

\textbf{Mean Average Precision (MAP)} — це середнє значення показника Average Precision для всіх користувачів.

Ця метрика враховує правильність рекомендацій та порядок, у якому вони показуються. Якщо релевантні об'єкти знаходяться на початку списку рекомендацій, значення MAP буде високим. Тому MAP широко використовується для оцінювання пошукових систем і рекомендаційних платформ.

\textbf{Normalized Discounted Cumulative Gain (NDCG)} оцінює якість ранжування рекомендованих об'єктів.

Особливість цієї метрики полягає в тому, що вона надає більшу вагу тим рекомендаціям, які знаходяться на початку списку. Таким чином, система отримує вищу оцінку, якщо найрелевантніші рекомендації показуються першими. NDCG особливо важлива для сервісів, де користувач переглядає лише перші кілька рекомендацій.

Під час оцінки якості рекомендаційної моделі часто використовується крос-валідація: дані розбиваються на декілька частин, і на кожній ітерації одна частина використовується для тестування, а решта — для навчання.


\ding{43} Сучасні рекомендаційні системи дедалі частіше використовують методи машинного навчання та глибокого навчання, зокрема градієнтний бустинг, нейронні мережі, автоенкодери (Autoencoders), рекурентні нейронні мережі (RNN), трансформери (Transformers), графові нейронні мережі (Graph Neural Networks, GNN) та моделі навчання з підкріпленням (Reinforcement Learning);


\section{Етичні проблеми та сучасні підходи.}

Попри численні переваги, рекомендаційні системи породжують низку етичних питань. 

Основними проблемами є порушення конфіденційності користувачів та збір великої кількості персональних даних. Сучасні рекомендаційні системи дедалі частіше проєктуються з урахуванням принципів етичного використання штучного інтелекту. 

Також варто пам'ятати за соціальну відповідальність. Сучасні рекомендаційні системи в соціальних платформах часто приводять до формування інформаційних бульбашок (Filter Bubble), що в свою чергу приводить до росту соціальної напруги між різними групами населення.

Останніми роками рекомендаційні системи активно інтегрують технології штучного інтелекту та великі мовні моделі (Large Language Models, LLM).

Поєднання рекомендаційних систем із сучасними моделями штучного інтелекту дозволяє створювати більш точні, гнучкі та персоналізовані сервіси, здатні враховувати як довгострокові вподобання, так і поточний контекст взаємодії користувача із системою.


\ding{46} \textbf{Питання для самоконтролю.}

\begin{enumerate}
\item Що таке рекомендаційна система?
\item Які основні завдання розв'язують рекомендаційні системи?
\item У яких сферах застосовуються рекомендаційні системи?
\item Які переваги використання рекомендаційних систем для користувачів і бізнесу?
\item Які основні типи рекомендаційних систем існують?
\item Що таке контентно-орієнтована фільтрація (Content-Based Filtering)?
\item На яких даних базуються контентно-орієнтовані рекомендації?
\item Які переваги та недоліки контентно-орієнтованого підходу?
\item Що таке колаборативна фільтрація (Collaborative Filtering)?
\item Чим User-Based Collaborative Filtering відрізняється від Item-Based Collaborative Filtering?
\item Як визначається схожість між користувачами?
\item Як визначається схожість між об'єктами?
\item Для чого використовується косинусна подібність у рекомендаційних системах?
\item Які інші міри схожості застосовуються в рекомендаційних системах?
\item Що таке матриця «користувач – об'єкт» (User–Item Matrix)?
\item Чому матриця оцінок зазвичай є розрідженою (Sparse Matrix)?
\item Що таке проблема «холодного старту» (Cold Start)?
\item Які існують способи розв'язання проблеми холодного старту?
\item Що таке гібридні рекомендаційні системи?
\item Які переваги мають гібридні рекомендаційні системи?
\item Для чого використовується матрична факторизація (Matrix Factorization)?
\item Як метод сингулярного розкладу (SVD) застосовується в рекомендаційних системах?
\item Які метрики використовуються для оцінювання рекомендаційних систем?
\item Що таке Precision та Recall у задачах рекомендації?
\item Що таке метрика Mean Average Precision (MAP)?
\item Для чого використовується метрика Normalized Discounted Cumulative Gain (NDCG)?
\item Як оцінюється якість рекомендацій за допомогою крос-валідації?
\item Які сучасні методи машинного навчання застосовуються для побудови рекомендаційних систем?
\item Які етичні проблеми можуть виникати під час використання рекомендаційних систем?
\item Які сучасні тенденції розвитку рекомендаційних систем пов'язані зі штучним інтелектом і великими мовними моделями (LLM)?
\end{enumerate}